\section{Oracoli e Meccanismi di
Verifica}\label{oracoli-e-meccanismi-di-verifica}

Per decidere se un team ha recuperato la \textbf{flag}, bisogna definire
un \textbf{\emph{oracolo}}, che determini lo stato in cui si trova il
sistema target.

L'oracolo è in grado di accedere al sistema e di fare due azioni: può
leggere la memoria fisica di qualsiasi nodo e può eseguire comandi su
qualsiasi nodo

Per ogni tipo di azione eseguita, l'oracolo determina lo stato del
sistema confrontando il risultato dell'azione con il valore che
l'oracolo si aspettava (ossia i valori vengono validati). In BlueAgent,
l'oracolo è l'entità che esegue sui container clonati i comandi
specificati dall'Autore della sfida, ossia nel codice la classe
\passthrough{\lstinline!OVNVerifier!}.

Dunque all'interno del processo di verifica, il sistema si può definire
\emph{patchato} se e soltanto se tutte le azioni sono validate rispetto
al valore che ci si aspettava in output.

Di seguito sono discussi alcuni meccanismi di verifica che abbiamo
identificato, al fine di motivare la clonazione dei container
implementata in BlueAgent.

\subsection{Verifica Statica}\label{verifica-statica}

La verifica statica prevede che l'oracolo possa accedere alla memoria
fisica, comparando le patch applicate nei file in memoria, ad una patch
considerata di riferimento per l'oracolo.

Questo tipo di verifica può essere implementato in vari modi:

\begin{itemize}
\tightlist
\item
  \textbf{diff}: ad ogni file interessato corrisponde un file di
  riferimento per l'oracolo. Ogni file viene dunque confrontato, riga
  per riga, con il suo riferimento per determinare se il team ha
  applicato la patch correttamente.
\item
  \textbf{SAST}: Static Analysis Security Testing. Prevede l'accesso al
  codice sorgente non compilato in modo da analizzare pattern di codice
  vulnerabili e fallacie logiche.\footnote{\label{fn:ref_13}https://aetic.theiaer.org/archive/v4/v4n3/p1.html
    2. Currrent Status of SAST based on Checkmarx}
\end{itemize}

Tuttavia, SAST non è un metodo di analisi efficace, sebbene efficiente:

\begin{itemize}
\tightlist
\item
  possono essere identificati nel processo di verifica, vari falsi
  positivi (e falsi negativi)\footref{fn:ref_13}
\item
  non posso identificare staticamente tutti i tipi di errori di
  sicurezza \footref{fn:ref_13}
\item
  non posso identificare vulnerabilità all'interno dei binari
  compilati\footref{fn:ref_13}
\end{itemize}

Dunque l'analisi statica può in alcuni casi risultare inesatta, e non
coprire tutti i casi d'uso per quanto riguarda la gestione di
competizioni CTF.

\subsection{Verifica Dinamica tramite sistema
esterno}\label{verifica-dinamica-tramite-sistema-esterno}

Questo meccanismo di verifica prevede l'uso di un \textbf{orchestratore
esterno} alle macchine vulnerabili, ossia l'attaccante, che possa
interagire con la vulnerabilità utilizzando la rete locale. L'attaccante
orchestra l'esecuzione di uno o più comandi, con l'obiettivo di inviare
il payload all'interno di un pacchetto di rete.

Durante una competizione di questo tipo, tuttavia, i membri dei team
sono legittimati ad accedere alla macchina vulnerabile come utente
\passthrough{\lstinline!root!}: così da permettere l'applicazione delle
patch sui file di sistema. Questo implica che i partecipanti dei team
siano in grado di eseguire tool come \passthrough{\lstinline!tcpdump!} e
\passthrough{\lstinline!wireshark!} per intercettare il traffico di rete
e di conseguenza, di visualizzarne il payload, ottenendo un suggerimento
per la soluzione della challenge.

Inoltre, il payload inviato potrebbe essere tale da mandare in crash il
servizio target, causando una penalità alla percentuale di \textbf{SLA}
\emph{(Service Level Agreement)}, ed interferendo potenzialmente con il
lavoro del team sfidante (portando per esempio alla perdita di alcuni
file). Si vuole dunque trovare un sistema di verifica che esegua
completamente all'oscuro del team sfidante.

Il sistema di verifica tramite sistema esterno non copre nemmeno tutte
le vulnerabilità. Banalmente, non tutti i servizi vulnerabili sono
esposti in rete, ma sono accessibili solo in locale.

Riassumendo quindi, abbiamo identificato alcuni spunti per ideare un
meccanismo di verifica migliore:

\begin{itemize}
\tightlist
\item
  Il traffico di rete non dovrebbe essere letto dal team sfidante al
  fine di ottenere suggerimenti per la soluzione.
\item
  Se l'exploit manda in crash il servizio, questo non dovrebbe
  penalizzare il team, o più in generale: il team sfidante non vede il
  servizio distruggersi.
\item
  Dobbiamo gestire anche vulnerabilità non esposte in rete.
\end{itemize}

\subsection{Verifica Dinamica tramite injection dei comandi di
verifica}\label{verifica-dinamica-tramite-injection-dei-comandi-di-verifica}

Questo meccanismo di verifica prevede l'uso di un \textbf{orchestratore
esterno} al sistema vulnerabile, che ha la facoltà di eseguire comandi
all'interno delle macchine vulnerabili che compongono il sistema. Dunque
l'orchestratore esegue uno o più comandi con il fine di verificare che
il sistema sia resistente alla vulnerabilità.

Considerando il caso in cui una macchina \passthrough{\lstinline!A!}
contenga un file binario vulnerabile ad un \emph{buffer overflow},
l'orchestratore potrebbe eseguire ad esempio il comando
\passthrough{\lstinline!/path/to/service\_logger payload!} e verificare
l'esito dell'esecuzione.

L'esecuzione di questi comandi è possibile per esempio in ambienti
containerizzati, sfruttando utility come
\passthrough{\lstinline!docker exec!} o
\passthrough{\lstinline!incus exec!}, che permettono di iniettare
comandi all'interno dei container.

Tuttavia: comandi come \passthrough{\lstinline!pspy!}, (che ascoltano
attivamente le chiamate di sistema effettuate e gli eventi sul file
system in \passthrough{\lstinline!/proc!}) possono intercettare
l'esecuzione dei comandi dell'orchestratore e di mostrarli al team,
suggerendo al team sfidante la soluzione. \footnote{\label{fn:ref_14}https://github.com/DominicBreuker/pspy}

Dunque, dobbiamo ideare un meccanismo di verifica che sia completamente
invisibile a questi meccanismi di intercettazione.

\subsection{Clonazione del sistema in una rete
separata}\label{clonazione-del-sistema-in-una-rete-separata}

Clonare i container all'interno di una rete separata ed inaccessibile ai
giocatori è una soluzione plausibile in quanto:

\begin{itemize}
\tightlist
\item
  Isola la rete rispetto alla rete dei giocatori. I giocatori non
  possono intercettare il traffico di rete.
\item
  I giocatori, per design, non hanno accesso ai container duplicati
  dunque non possono eseguire alcun tipo di comando.
\end{itemize}

Tuttavia introduce un costo in termini di storage e di utilizzo della
RAM dovuto alla duplicazione dei container, anche se questi esistono
solo per il tempo necessario all'esecuzione dei comandi definiti dal
creatore della sfida. Questo aspetto verrà misurato più avanti e sarà
oggetto di analisi.

Altri meccanismi di verifica ed architetture di rete possibili sono
analizzati nel Capitolo 6.
